%%% KLOCKI: TABELE WYNIKÓW
%%% Wiersze i kolumny wynikają z ustawień eksperymentów. Wyniki wpisujesz w miejsce \wyn.
%%% Do pracy tylko wyniki z aplikacji IL2CPP w wersji Release (kolumna srodowisko równa IL2CPP).
%%% Pliki z folderu Wyniki: główny E1_*.csv oraz podsumowania E1_*_warianty.csv i E1_*_korelacje.csv.

\CoTuMaByc{Z \kol{E1_*_warianty.csv}: \kol{czas_ms_mediana}, \kol{czas_ms_q1}, \kol{czas_ms_q3}, \kol{pamiec_rzeczywista_mb_mediana}. Z głównego pliku: \kol{czas_propagacja_ms}.}

\begin{table}
  \centering
  \footnotesize
  \caption{E1: czas generacji i pamięć według rozwiązywacza}
  \label{tab:e1}
  \begin{tabularx}{\linewidth}{@{}lCCCCC@{}}
    \toprule
    Rozwiązywacz & Czas: mediana [ms] & Czas: Q1 [ms] & Czas: Q3 [ms] & Propagacja: mediana [ms] & Pamięć rzeczywista [MB] \\
    \midrule
    \kol{Naive} & 837,8 & 803,4 & 1449,1 & 825,2 & 6,12 \\
    \kol{SupportCount} & 81,3 & 79,9 & 132,0 & 67,3 & 13,43 \\
    \kol{Bitset} & 26,3 & 24,7 & 36,4 & 12,1 & 1,92 \\
    \kol{Burst} & 20,1 & 18,4 & 25,3 & 5,8 & 1,92 \\
    \bottomrule
  \end{tabularx}
\end{table}

\CoTuMaByc{Z \kol{*_mann_whitney.csv} ze skryptu \kol{Analiza/analiza_wynikow.py}: \kol{U}, \kol{p}. Test: \cite{MannWhitney1947}.}

\begin{table}
  \centering
  \footnotesize
  \caption{E1: test Manna-Whitneya dla par rozwiązywaczy}
  \label{tab:e1-mann-whitney}
  \begin{tabular}{@{}llrr@{}}
    \toprule
    Wariant A & Wariant B & $U$ & $p$ \\
    \midrule
    \kol{Naive} & \kol{SupportCount} & 10\,000 & $2{,}56 \cdot 10^{-34}$ \\
    \kol{Naive} & \kol{Bitset} & 10\,000 & $2{,}56 \cdot 10^{-34}$ \\
    \kol{Naive} & \kol{Burst} & 10\,000 & $2{,}56 \cdot 10^{-34}$ \\
    \kol{SupportCount} & \kol{Bitset} & 9928 & $2{,}19 \cdot 10^{-33}$ \\
    \kol{SupportCount} & \kol{Burst} & 10\,000 & $2{,}56 \cdot 10^{-34}$ \\
    \kol{Bitset} & \kol{Burst} & 8107 & $3{,}19 \cdot 10^{-14}$ \\
    \bottomrule
  \end{tabular}
\end{table}

\CoTuMaByc{Z \kol{E2_*_warianty.csv}: \kol{czas_ms_mediana}, \kol{proby}, \kol{sprzecznosci}, \kol{udane_proby_proc}.}

\begin{table}
  \centering
  \footnotesize
  \caption[E2: czas i sprzeczności według heurystyki]{E2: czas i sprzeczności według heurystyki. Próby i sprzeczności: suma ze 100 map}
  \label{tab:e2}
  \begin{tabularx}{\linewidth}{@{}lCCCC@{}}
    \toprule
    Heurystyka & Czas: mediana [ms] & Próby & Sprzeczności & Udane próby [\%] \\
    \midrule
    \kol{LevelFirst} & 26,4 & 139 & 39 & 71,9 \\
    \kol{Scanline} & 16,3 & 159 & 59 & 62,9 \\
    \kol{MinRemaining} & 226,7 & 2649 & 2549 & 3,8 \\
    \kol{Entropy} & 330,7 & 4086 & 3986 & 2,4 \\
    \kol{Random} & 381,7 & 5765 & 5665 & 1,7 \\
    \bottomrule
  \end{tabularx}
\end{table}

\CoTuMaByc{Z głównego pliku E2: \kol{wieze}, \kol{donzony}, \kol{mosty}, \kol{schody}, \kol{wys_wiez_sr}, \kol{wys_murow_sr}, \kol{trawa_dostepna_proc}.}

\begin{table}
  \centering
  \footnotesize
  \caption[E2: budowa map według heurystyki]{E2: budowa map według heurystyki, średnie na mapę}
  \label{tab:e2-budowa}
  \begin{tabularx}{\linewidth}{@{}lCCCCCCC@{}}
    \toprule
    Heurystyka & Wieże & Donżony & Mosty & Schody & Wysokość wież & Wysokość murów & Teren osiągalny [\%] \\
    \midrule
    \kol{LevelFirst} & 11,9 & 6,5 & 1,2 & 4,4 & 5,7 & 2,9 & 97,1 \\
    \kol{Scanline} & 16,6 & 4,5 & 1,5 & 3,7 & 5,9 & 2,9 & 94,1 \\
    \kol{MinRemaining} & 26,8 & 2,5 & 35,2 & 15,1 & 11,7 & 8,4 & 42,8 \\
    \kol{Entropy} & 29,2 & 1,7 & 26,6 & 10,9 & 11,6 & 9,0 & 41,6 \\
    \kol{Random} & 22,4 & 7,4 & 12,8 & 12,2 & 10,5 & 7,0 & 72,9 \\
    \bottomrule
  \end{tabularx}
\end{table}

\CoTuMaByc{Odsetek komórek innych niż w mapie poprzedniego ziarna. Z głównego pliku E2: \kol{roznica_mapy_proc}, \kol{roznica_parter_proc}, \kol{roznica_zabudowa_proc}.}

\begin{table}
  \centering
  \footnotesize
  \caption[E2: różnorodność map według heurystyki]{E2: różnorodność map według heurystyki, średni odsetek komórek innych niż w mapie poprzedniego ziarna}
  \label{tab:roznorodnosc}
  \begin{tabularx}{\linewidth}{@{}lCCC@{}}
    \toprule
    Heurystyka & Cała mapa [\%] & Parter [\%] & Zabudowa [\%] \\
    \midrule
    \kol{LevelFirst} & 23,9 & 79,5 & 98,2 \\
    \kol{Scanline} & 25,4 & 80,5 & 97,9 \\
    \kol{MinRemaining} & 67,8 & 82,7 & 94,6 \\
    \kol{Entropy} & 69,2 & 81,4 & 93,7 \\
    \kol{Random} & 58,5 & 83,2 & 96,2 \\
    \bottomrule
  \end{tabularx}
\end{table}

\CoTuMaByc{Z \kol{E3_*_warianty.csv}: \kol{czas_ms_mediana}. Limit 600 s na mapę; przekroczenie zapisuje się jako wynik.}

\begin{table}
  \centering
  \footnotesize
  \caption{E3: mediana czasu generacji [ms]}
  \label{tab:e3-czas}
  \begin{tabular}{@{}lrrrr@{}}
    \toprule
    Reakcja & 10$\times$10$\times$8 & 20$\times$20$\times$13 & 40$\times$40$\times$13 & 80$\times$80$\times$13 \\
    \midrule
    \kol{Restart} & 3,0 & 26,3 & 300,1 & 114\,750 \\
    \kol{Backtrack} & 3,0 & 25,7 & 117,1 & 496,6 \\
    \kol{LocalReset} & 3,0 & 25,4 & 116,7 & 497,2 \\
    \bottomrule
  \end{tabular}
\end{table}

\CoTuMaByc{Z \kol{E3_*_warianty.csv}: \kol{udane_mapy}, \kol{mapy}.}

\begin{table}
  \centering
  \footnotesize
  \caption{E3: udane mapy na liczbę map}
  \label{tab:e3-powodzenie}
  \begin{tabular}{@{}lrrrr@{}}
    \toprule
    Reakcja & 10$\times$10$\times$8 & 20$\times$20$\times$13 & 40$\times$40$\times$13 & 80$\times$80$\times$13 \\
    \midrule
    \kol{Restart} & 100/100 & 100/100 & 100/100 & 50/100 \\
    \kol{Backtrack} & 100/100 & 100/100 & 100/100 & 100/100 \\
    \kol{LocalReset} & 100/100 & 100/100 & 100/100 & 100/100 \\
    \bottomrule
  \end{tabular}
\end{table}

\CoTuMaByc{Z \kol{E4_*_warianty.csv}: \kol{czas_ms_mediana}. Nachylenie: \cref{eq:nachylenie}. Limit 600 s na mapę.}

\begin{table}
  \centering
  \footnotesize
  \caption[E4: mediana czasu generacji]{E4: mediana czasu generacji [ms]. Nachylenie $a$ dopasowano do trzech mniejszych siatek, bo przy 80$\times$80$\times$13 część map nie powstaje}
  \label{tab:e4-czas}
  \begin{tabular}{@{}lrrrrr@{}}
    \toprule
    Rozwiązywacz & 10$\times$10$\times$8 & 20$\times$20$\times$13 & 40$\times$40$\times$13 & 80$\times$80$\times$13 & Nachylenie $a$ \\
    \midrule
    Komórki & 800 & 5200 & 20\,800 & 83\,200 & -- \\
    \midrule
    \kol{Naive} & 90,0 & 835,5 & 12\,666 & 601\,188 & 1,50 \\
    \kol{SupportCount} & 10,3 & 82,9 & 1229,5 & 600\,000 & 1,45 \\
    \kol{Bitset} & 3,1 & 26,0 & 297,2 & 115\,059 & 1,39 \\
    \kol{Burst} & 2,3 & 19,8 & 207,3 & 73\,773 & 1,36 \\
    \bottomrule
  \end{tabular}
\end{table}

\CoTuMaByc{Z \kol{E4_*_warianty.csv}: \kol{pamiec_rzeczywista_mb_mediana}. Pojedyncze wartości pamięci szumią, więc porównuj mediany.}

\begin{table}
  \centering
  \footnotesize
  \caption{E4: mediana pamięci rzeczywistej [MB]}
  \label{tab:e4-pamiec}
  \begin{tabular}{@{}lrrrrr@{}}
    \toprule
    Rozwiązywacz & 10$\times$10$\times$8 & 20$\times$20$\times$13 & 40$\times$40$\times$13 & 80$\times$80$\times$13 & Nachylenie $a$ \\
    \midrule
    Komórki & 800 & 5200 & 20\,800 & 83\,200 & -- \\
    \midrule
    \kol{Naive} & 0,88 & 6,10 & 24,41 & 85,72 & 0,99 \\
    \kol{SupportCount} & 2,15 & 13,43 & 53,23 & 212,45 & 0,99 \\
    \kol{Bitset} & 0,26 & 1,92 & 7,57 & 30,22 & 1,02 \\
    \kol{Burst} & 0,26 & 1,92 & 7,57 & 30,22 & 1,02 \\
    \bottomrule
  \end{tabular}
\end{table}

\CoTuMaByc{Z \kol{E5_*_warianty.csv}: \kol{czas_ms_mediana}, \kol{czas_ms_q1}, \kol{czas_ms_q3}, \kol{sprzecznosci}.}

\begin{table}
  \centering
  \footnotesize
  \caption[E5: czas generacji według zestawu]{E5: czas generacji według zestawu. Sprzeczności: suma ze 100 map}
  \label{tab:e5}
  \setlength{\tabcolsep}{3pt}
  \begin{tabularx}{\linewidth}{@{}lCCCCCC@{}}
    \toprule
    Zestaw & Kafelki & Stany & Czas: mediana [ms] & Czas: Q1 [ms] & Czas: Q3 [ms] & Sprzeczności \\
    \midrule
    \kol{Prosty} & 8 & 26 & 5,8 & 4,4 & 8,2 & 117 \\
    \kol{Sredni} & 22 & 79 & 11,7 & 11,4 & 12,6 & 26 \\
    \kol{PelnyBezDekoracji} & 41 & 146 & 25,9 & 24,7 & 36,6 & 60 \\
    \kol{Pelny} & 50 & 170 & 26,0 & 24,7 & 36,4 & 39 \\
    \bottomrule
  \end{tabularx}
\end{table}

\CoTuMaByc{Z \kol{E6_*_warianty.csv}: \kol{sprzecznosci}, \kol{udane_proby_proc}, \kol{zabudowa_proc}. Z głównego pliku: \kol{wieze}.}

\begin{table}
  \centering
  \footnotesize
  \caption[E6: wyniki dla zestawu Pelny i wariantów reguł]{E6: wyniki dla zestawu Pelny i wariantów reguł. Sprzeczności: suma ze 100 map, wieże: średnia na mapę}
  \label{tab:e6}
  \begin{tabularx}{\linewidth}{@{}lCCCC@{}}
    \toprule
    Wariant & Sprzeczności & Udane próby [\%] & Zabudowa [\%] & Wieże \\
    \midrule
    \kol{Pelny} & 39 & 71,9 & 39,4 & 11,9 \\
    \kol{PelnyBezWolnegoPola} & 958 & 9,5 & 56,6 & 24,2 \\
    \kol{PelnyBezDrzwiNadSchodami} & 13 & 88,5 & 40,4 & 12,6 \\
    \kol{PelnyTrawa040} & 113 & 46,9 & 50,6 & 15,9 \\
    \kol{PelnyTrawa025} & 210 & 32,3 & 55,1 & 16,4 \\
    \bottomrule
  \end{tabularx}
\end{table}

\CoTuMaByc{Z \kol{E7_*_warianty.csv} albo mediana z głównego pliku: \kol{wywolania_rysowania_mediana}.}

\begin{table}
  \centering
  \footnotesize
  \caption{E7: mediana wywołań rysowania na klatkę}
  \label{tab:e7-wywolania}
  \begin{tabular}{@{}llrrrr@{}}
    \toprule
    Rysowanie & Łączenie rysowania & 10$\times$10$\times$8 & 20$\times$20$\times$13 & 40$\times$40$\times$13 & 80$\times$80$\times$13 \\
    \midrule
    \multirow{3}{*}{\kol{GameObjects}} & SRP Batcher + Resident Drawer & 115,5 & 183 & 211 & 211 \\
     & Sam SRP Batcher & 692 & 2916 & 12\,606 & 51\,603 \\
     & Samo instancjonowanie & 148 & 304 & 417 & 507 \\
    \midrule
    \multirow{3}{*}{\kol{Instanced}} & SRP Batcher + Resident Drawer & 216,5 & 432 & 213 & 267,5 \\
     & Sam SRP Batcher & 217 & 432 & 213 & 267,5 \\
     & Samo instancjonowanie & 217 & 432 & 213 & 267,5 \\
    \midrule
    \multirow{3}{*}{\kol{Combined}} & SRP Batcher + Resident Drawer & 39 & 122 & 225 & 613 \\
     & Sam SRP Batcher & 39 & 122 & 225 & 613 \\
     & Samo instancjonowanie & 39 & 122 & 225 & 613 \\
    \midrule
    \multirow{3}{*}{\kol{CombinedAtlas}} & SRP Batcher + Resident Drawer & 19 & 27 & 35 & 67 \\
     & Sam SRP Batcher & 19 & 27 & 35 & 67 \\
     & Samo instancjonowanie & 19 & 27 & 35 & 67 \\
    \bottomrule
  \end{tabular}
\end{table}

\CoTuMaByc{Z \kol{E7_*_warianty.csv} albo mediana z głównego pliku: \kol{narysowane_obiekty_mediana}.}

\begin{table}
  \centering
  \footnotesize
  \caption{E7: mediana narysowanych obiektów na wywołanie rysowania}
  \label{tab:e7-obiekty}
  \begin{tabular}{@{}llrrrr@{}}
    \toprule
    Rysowanie & Łączenie rysowania & 10$\times$10$\times$8 & 20$\times$20$\times$13 & 40$\times$40$\times$13 & 80$\times$80$\times$13 \\
    \midrule
    \multirow{3}{*}{\kol{GameObjects}} & SRP Batcher + Resident Drawer & 5,8 & 15,9 & 60,4 & 244,6 \\
     & Sam SRP Batcher & 1,0 & 1,0 & 1,0 & 1,0 \\
     & Samo instancjonowanie & 4,7 & 9,5 & 30,3 & 101,8 \\
    \midrule
    \multirow{3}{*}{\kol{Instanced}} & SRP Batcher + Resident Drawer & 6,1 & 16,7 & 59,8 & 194,0 \\
     & Sam SRP Batcher & 6,1 & 16,7 & 59,8 & 194,0 \\
     & Samo instancjonowanie & 6,1 & 16,7 & 59,8 & 194,0 \\
    \midrule
    \multirow{3}{*}{\kol{Combined}} & SRP Batcher + Resident Drawer & 1,0 & 1,0 & 1,0 & 1,0 \\
     & Sam SRP Batcher & 1,0 & 1,0 & 1,0 & 1,0 \\
     & Samo instancjonowanie & 1,0 & 1,0 & 1,0 & 1,0 \\
    \midrule
    \multirow{3}{*}{\kol{CombinedAtlas}} & SRP Batcher + Resident Drawer & 1,0 & 1,0 & 1,0 & 1,0 \\
     & Sam SRP Batcher & 1,0 & 1,0 & 1,0 & 1,0 \\
     & Samo instancjonowanie & 1,0 & 1,0 & 1,0 & 1,0 \\
    \bottomrule
  \end{tabular}
\end{table}

\CoTuMaByc{Z \kol{E7_*_warianty.csv} albo mediana z głównego pliku: \kol{setpass}.}

\begin{table}
  \centering
  \footnotesize
  \caption{E7: mediana zmian stanu SetPass na klatkę}
  \label{tab:e7-setpass}
  \begin{tabular}{@{}llrrrr@{}}
    \toprule
    Rysowanie & Łączenie rysowania & 10$\times$10$\times$8 & 20$\times$20$\times$13 & 40$\times$40$\times$13 & 80$\times$80$\times$13 \\
    \midrule
    \multirow{3}{*}{\kol{GameObjects}} & SRP Batcher + Resident Drawer & 9 & 9 & 9 & 9 \\
     & Sam SRP Batcher & 9 & 15 & 37 & 37 \\
     & Samo instancjonowanie & 45 & 61 & 83 & 83 \\
    \midrule
    \multirow{3}{*}{\kol{Instanced}} & SRP Batcher + Resident Drawer & 31 & 34 & 31 & 31 \\
     & Sam SRP Batcher & 31 & 34 & 31 & 31 \\
     & Samo instancjonowanie & 31 & 34 & 31 & 31 \\
    \midrule
    \multirow{3}{*}{\kol{Combined}} & SRP Batcher + Resident Drawer & 9 & 11 & 9 & 11 \\
     & Sam SRP Batcher & 9 & 11 & 9 & 9 \\
     & Samo instancjonowanie & 29 & 55 & 55 & 55 \\
    \midrule
    \multirow{3}{*}{\kol{CombinedAtlas}} & SRP Batcher + Resident Drawer & 9 & 11 & 9 & 9 \\
     & Sam SRP Batcher & 9 & 11 & 9 & 9 \\
     & Samo instancjonowanie & 9 & 11 & 9 & 9 \\
    \bottomrule
  \end{tabular}
\end{table}

\CoTuMaByc{Z \kol{E7_*_warianty.csv} albo mediana z głównego pliku: \kol{czas_klatki_ms_mediana}.}

\begin{table}
  \centering
  \footnotesize
  \caption{E7: mediana czasu klatki [ms]}
  \label{tab:e7-klatka}
  \begin{tabular}{@{}llrrrr@{}}
    \toprule
    Rysowanie & Łączenie rysowania & 10$\times$10$\times$8 & 20$\times$20$\times$13 & 40$\times$40$\times$13 & 80$\times$80$\times$13 \\
    \midrule
    \multirow{3}{*}{\kol{GameObjects}} & SRP Batcher + Resident Drawer & 0,513 & 0,478 & 0,552 & 0,691 \\
     & Sam SRP Batcher & 0,487 & 0,717 & 1,707 & 6,039 \\
     & Samo instancjonowanie & 0,509 & 0,747 & 1,446 & 4,252 \\
    \midrule
    \multirow{3}{*}{\kol{Instanced}} & SRP Batcher + Resident Drawer & 0,558 & 0,630 & 0,683 & 1,126 \\
     & Sam SRP Batcher & 0,520 & 0,625 & 0,624 & 1,046 \\
     & Samo instancjonowanie & 0,513 & 0,638 & 0,651 & 1,091 \\
    \midrule
    \multirow{3}{*}{\kol{Combined}} & SRP Batcher + Resident Drawer & 0,426 & 0,505 & 0,468 & 0,560 \\
     & Sam SRP Batcher & 0,370 & 0,411 & 0,457 & 0,508 \\
     & Samo instancjonowanie & 0,393 & 0,484 & 0,535 & 0,690 \\
    \midrule
    \multirow{3}{*}{\kol{CombinedAtlas}} & SRP Batcher + Resident Drawer & 0,432 & 0,444 & 0,453 & 0,514 \\
     & Sam SRP Batcher & 0,365 & 0,405 & 0,475 & 0,458 \\
     & Samo instancjonowanie & 0,364 & 0,415 & 0,469 & 0,461 \\
    \bottomrule
  \end{tabular}
\end{table}

\CoTuMaByc{Z \kol{E7_*_warianty.csv} albo mediana z głównego pliku: \kol{najdluzsza_klatka_ms_mediana}.}

\begin{table}
  \centering
  \footnotesize
  \caption{E7: mediana najdłuższej klatki [ms]}
  \label{tab:e7-najdluzsza}
  \begin{tabular}{@{}llrrrr@{}}
    \toprule
    Rysowanie & Łączenie rysowania & 10$\times$10$\times$8 & 20$\times$20$\times$13 & 40$\times$40$\times$13 & 80$\times$80$\times$13 \\
    \midrule
    \multirow{3}{*}{\kol{GameObjects}} & SRP Batcher + Resident Drawer & 1,38 & 1,20 & 1,35 & 1,80 \\
     & Sam SRP Batcher & 1,20 & 1,04 & 2,17 & 8,08 \\
     & Samo instancjonowanie & 1,20 & 1,14 & 1,90 & 6,26 \\
    \midrule
    \multirow{3}{*}{\kol{Instanced}} & SRP Batcher + Resident Drawer & 1,27 & 1,23 & 1,32 & 2,03 \\
     & Sam SRP Batcher & 1,21 & 1,19 & 1,32 & 1,93 \\
     & Samo instancjonowanie & 1,21 & 1,20 & 1,31 & 2,12 \\
    \midrule
    \multirow{3}{*}{\kol{Combined}} & SRP Batcher + Resident Drawer & 1,13 & 1,09 & 1,23 & 1,16 \\
     & Sam SRP Batcher & 1,15 & 1,21 & 1,23 & 1,18 \\
     & Samo instancjonowanie & 1,17 & 1,23 & 1,25 & 1,20 \\
    \midrule
    \multirow{3}{*}{\kol{CombinedAtlas}} & SRP Batcher + Resident Drawer & 1,14 & 1,18 & 1,21 & 1,27 \\
     & Sam SRP Batcher & 1,14 & 1,16 & 1,38 & 1,28 \\
     & Samo instancjonowanie & 1,17 & 1,15 & 1,27 & 1,31 \\
    \bottomrule
  \end{tabular}
\end{table}

\CoTuMaByc{Z \kol{E7_*_warianty.csv}: \kol{czas_stawiania_ms_mediana}. Podaj, dla którego ustawienia łączenia rysowania.}

\begin{table}
  \centering
  \footnotesize
  \caption[E7: mediana czasu stawiania mapy]{E7: mediana czasu stawiania mapy [ms], sam SRP Batcher}
  \label{tab:e7-stawianie}
  \begin{tabular}{@{}lrrrr@{}}
    \toprule
    Rysowanie & 10$\times$10$\times$8 & 20$\times$20$\times$13 & 40$\times$40$\times$13 & 80$\times$80$\times$13 \\
    \midrule
    \kol{GameObjects} & 0,48 & 2,08 & 10,98 & 62,10 \\
    \kol{Instanced} & 0,15 & 0,33 & 0,85 & 3,25 \\
    \kol{Combined} & 0,47 & 1,15 & 3,53 & 11,14 \\
    \kol{CombinedAtlas} & 0,52 & 1,07 & 3,28 & 10,38 \\
    \bottomrule
  \end{tabular}
\end{table}

\CoTuMaByc{Z \kol{E7_*_warianty.csv}: \kol{pamiec_unity_mb_mediana}. Podaj, dla którego ustawienia łączenia rysowania.}

\begin{table}
  \centering
  \footnotesize
  \caption[E7: mediana pamięci silnika]{E7: mediana pamięci silnika [MB], sam SRP Batcher}
  \label{tab:e7-pamiec}
  \begin{tabular}{@{}lrrrr@{}}
    \toprule
    Rysowanie & 10$\times$10$\times$8 & 20$\times$20$\times$13 & 40$\times$40$\times$13 & 80$\times$80$\times$13 \\
    \midrule
    \kol{GameObjects} & 303,8 & 313,1 & 358,7 & 397,9 \\
    \kol{Instanced} & 302,6 & 307,7 & 353,3 & 368,7 \\
    \kol{Combined} & 297,0 & 308,8 & 370,1 & 425,7 \\
    \kol{CombinedAtlas} & 297,0 & 308,7 & 370,1 & 425,9 \\
    \bottomrule
  \end{tabular}
\end{table}

\CoTuMaByc{Z głównego pliku E7: \kol{pamiec_siatek_mb}, \kol{bok_fragmentu}.}

\begin{table}
  \centering
  \footnotesize
  \caption[E7: rozmiar scalonych siatek]{E7: mediana rozmiaru scalonych siatek [MB], fragmenty o boku 16 komórek}
  \label{tab:e7-siatki}
  \begin{tabular}{@{}lrrrr@{}}
    \toprule
    Rysowanie & 10$\times$10$\times$8 & 20$\times$20$\times$13 & 40$\times$40$\times$13 & 80$\times$80$\times$13 \\
    \midrule
    \kol{Combined} & 0,55 & 2,31 & 10,57 & 44,96 \\
    \kol{CombinedAtlas} & 0,55 & 2,31 & 10,57 & 44,96 \\
    \bottomrule
  \end{tabular}
\end{table}

\CoTuMaByc{Z \kol{E8_*_warianty.csv}: \kol{gestosc_regul_proc}, \kol{gestosc_wazona_proc}, \kol{kappa}, \kol{udane_proby_proc}, \kol{zabudowa_proc}.}

\begin{table}
  \centering
  \footnotesize
  \caption{E8: restrykcyjność reguł i powodzenie generacji}
  \label{tab:e8}
  \begin{tabularx}{\linewidth}{@{}lCCCCCC@{}}
    \toprule
    Zestaw lub wariant & Stany & Gęstość par [\%] & Gęstość ważona [\%] & $\kappa$ & Udane próby [\%] & Zabudowa [\%] \\
    \midrule
    \kol{Prosty} & 26 & 7,79 & 21,62 & 2,350 & 46,1 & 29,1 \\
    \kol{Sredni} & 79 & 10,24 & 18,99 & 1,565 & 79,4 & 35,8 \\
    \kol{PelnyBezDekoracji} & 146 & 8,45 & 12,10 & 1,487 & 62,5 & 40,5 \\
    \kol{Pelny} & 170 & 8,05 & 11,67 & 1,471 & 71,9 & 39,4 \\
    \midrule
    \kol{PelnyBezWolnegoPola} & 170 & 13,37 & 15,65 & 1,175 & 9,5 & 56,6 \\
    \kol{PelnyBezDrzwiNadSchodami} & 170 & 8,17 & 11,79 & 1,463 & 88,5 & 40,4 \\
    \kol{PelnyTrawa040} & 170 & 8,05 & 11,06 & 1,471 & 46,9 & 50,6 \\
    \kol{PelnyTrawa025} & 170 & 8,05 & 10,93 & 1,471 & 32,3 & 55,1 \\
    \bottomrule
  \end{tabularx}
\end{table}

\CoTuMaByc{Z \kol{*_korelacje.csv}: \kol{miara}, \kol{warianty}, \kol{liczba_wariantow}, \kol{spearman_z_udanymi_probami}. Mapa zdegenerowana w kodzie: zabudowa parteru poniżej 1\% albo powyżej 90\%. Korelacja rang: \cite{Spearman1904}.}

\begin{table}
  \centering
  \footnotesize
  \caption{Korelacja Spearmana restrykcyjności z odsetkiem udanych prób}
  \label{tab:korelacje}
  \begin{tabular}{@{}llrrrr@{}}
    \toprule
    Miara & Warianty & E8: liczba & E8: $\rho$ & E9: liczba & E9: $\rho$ \\
    \midrule
    \multirow{2}{*}{Gęstość par} & wszystkie & 8 & 0,098 & 81 & 0,388 \\
     & bez map zdegenerowanych & 8 & 0,098 & 75 & 0,531 \\
    \midrule
    \multirow{2}{*}{Gęstość ważona} & wszystkie & 8 & 0,048 & 81 & 0,430 \\
     & bez map zdegenerowanych & 8 & 0,048 & 75 & 0,569 \\
    \midrule
    \multirow{2}{*}{$\kappa$} & wszystkie & 8 & 0,146 & 81 & $-$0,388 \\
     & bez map zdegenerowanych & 8 & 0,146 & 75 & $-$0,531 \\
    \bottomrule
  \end{tabular}
\end{table}

\CoTuMaByc{Z \kol{E9_*_warianty.csv}: \kol{zaostrzenie_proc}, \kol{losowanie_zaostrzenia}, \kol{udane_proby_proc}, \kol{zabudowa_proc}.}

\begin{table}
  \centering
  \footnotesize
  \caption{E9: wyniki według poziomu zaostrzenia reguł}
  \label{tab:e9}
  \begin{tabularx}{\linewidth}{@{}CCCCCCC@{}}
    \toprule
    Zaostrzenie [\%] & Losowania & Udane próby: mediana [\%] & Udane próby: minimum [\%] & Udane próby: maksimum [\%] & Zabudowa: mediana [\%] & Warianty z mapami zdegenerowanymi \\
    \midrule
    0 & 1 & 71,4 & 71,4 & 71,4 & 39,6 & 0 \\
    0,2 & 10 & 64,5 & 58,8 & 74,1 & 39,9 & 0 \\
    0,5 & 10 & 31,7 & 0,8 & 66,7 & 40,1 & 0 \\
    1 & 10 & 30,0 & 0,4 & 66,7 & 38,5 & 1 \\
    2 & 10 & 7,9 & 0,0 & 76,9 & 35,6 & 1 \\
    3 & 10 & 6,0 & 0,0 & 60,6 & 36,9 & 1 \\
    5 & 10 & 17,7 & 0,0 & 87,0 & 26,7 & 0 \\
    8 & 10 & 11,3 & 0,0 & 100,0 & 3,7 & 1 \\
    12 & 10 & 1,0 & 0,0 & 100,0 & 0,0 & 2 \\
    \bottomrule
  \end{tabularx}
\end{table}

\CoTuMaByc{Z głównego pliku E1: \kol{decyzje}, \kol{kroki_propagacji}, \kol{usuniete_propagacja}, \kol{usuniete_decyzje}. Udział i usunięcia na decyzję: \cref{eq:propagacja}.}

\begin{table}
  \centering
  \footnotesize
  \caption[Propagacja według rozwiązywacza (E1)]{Propagacja według rozwiązywacza (E1), mediany na mapę}
  \label{tab:propagacja-rozwiazywacze}
  \begin{tabularx}{\linewidth}{@{}lCCCCCC@{}}
    \toprule
    Rozwiązywacz & Decyzje & Kroki propagacji & Usunięte przez propagację & Usunięte przez decyzje & Udział propagacji [\%] & Usunięcia na decyzję \\
    \midrule
    \kol{Naive} & 702 & 273\,449 & 684\,265 & 2564 & 99,66 & 1092,3 \\
    \kol{SupportCount} & 702 & 4\,030\,944 & 684\,265 & 2564 & 99,66 & 1092,3 \\
    \kol{Bitset} & 702 & 273\,449 & 684\,265 & 2564 & 99,66 & 1092,3 \\
    \kol{Burst} & 702 & 273\,449 & 684\,265 & 2564 & 99,66 & 1092,3 \\
    \bottomrule
  \end{tabularx}
\end{table}

\CoTuMaByc{Z głównego pliku E4; podaj, dla którego rozwiązywacza.}

\begin{table}
  \centering
  \footnotesize
  \caption[Propagacja według rozmiaru siatki (E4)]{Propagacja według rozmiaru siatki (E4), rozwiązywacz Bitset, mediany na mapę. Liczniki obejmują wszystkie próby, razem z restartami}
  \label{tab:propagacja-rozmiary}
  \begin{tabularx}{\linewidth}{@{}lCCCCCC@{}}
    \toprule
    Siatka & Decyzje & Kroki propagacji & Usunięte przez propagację & Usunięte przez decyzje & Udział propagacji [\%] & Usunięcia na decyzję \\
    \midrule
    10$\times$10$\times$8 & 160 & 35\,891 & 84\,153 & 601 & 99,33 & 566,5 \\
    20$\times$20$\times$13 & 702 & 273\,449 & 684\,265 & 2564 & 99,66 & 1092,3 \\
    40$\times$40$\times$13 & 5504 & 3\,436\,857 & 10\,646\,474 & 25\,646 & 99,76 & 2011,5 \\
    80$\times$80$\times$13 & 991\,951 & 1\,233\,864\,193 & 5\,603\,406\,919 & 5\,882\,402 & 99,89 & 5490,4 \\
    \bottomrule
  \end{tabularx}
\end{table}

\CoTuMaByc{Kolumny głównych plików: \kol{udane}, \kol{naruszenia}.}

\begin{table}
  \centering
  \footnotesize
  \caption{Poprawność map we wszystkich eksperymentach}
  \label{tab:naruszenia}
  \begin{tabular}{@{}lrrr@{}}
    \toprule
    Eksperyment & Uruchomienia & Udane mapy & Naruszenia zasad \\
    \midrule
    E1 & 400 & 400 & 0 \\
    E2 & 500 & 500 & 0 \\
    E3 & 1200 & 1150 & 0 \\
    E4 & 1600 & 1350 & 0 \\
    E5 & 400 & 400 & 0 \\
    E6 & 500 & 500 & 0 \\
    E7 & 1440 & 1440 & 0 \\
    E8 & 800 & 800 & 0 \\
    E9 & 1620 & 1100 & 0 \\
    \bottomrule
  \end{tabular}
\end{table}
